% R15-3 -- Soundness/reduction/termination lemmas for the computer-assisted
% part of the C_3 program, ready to integrate into Section 10 (the
% covering-constant program) when the numbering freeze lifts.  Pure
% mathematics; the SHA/counts of Pillar 5 go into Appendix B separately.
% Labels are placeholders (thm:extreme etc.) to be reconciled on integration.

\subsection*{A partially computer-assisted case of the program}

Fix the cyclic triple $\tau_0=(\tfrac{13}{25},\tfrac12,\tfrac12)$ with
directions $u_1,u_2,u_3$ (values $(\tau_0)$ at $A,B,C$ cyclically), which is
\emph{non-concurrent} ($\Pi=\tfrac{13}{100}\ne\tfrac{12}{100}=Q$, so
$\delta>0$ by Theorem~\ref{thm:canon}); it is the ``sandwich'' of
Remark~\ref{rem:sandwich}. We prove $C_\Delta(\tau_0)=1$: the
covering-constant program (Problem~\ref{prob:covconst}) holds at $\tau_0$,
the first non-concurrent triple for which we establish it. A covering here
means three planks, one per direction, $P_i=\{u_i\in[l_i,h_i]\}$ with
$0\le l_i\le h_i\le1$ and $r_i=h_i-l_i$.

The proof reduces, by four elementary lemmas, to the emptiness of an
explicit finite bisection tree, which is verified by an exact rational
computation and re-verified by an independent checker (Appendix~\ref{app:cert}).

\begin{lemma}[extreme regime]\label{lem:extreme}
Set $w_0(\tau)=\tfrac12\min$ over the three cyclic rotations of
$\min\bigl(\tau_1,\,1-\tau_1,\,\tfrac{(1-\tau_2)\tau_1}{1+\tau_1},\,
\tfrac{\tau_3(1-\tau_1)}{2-\tau_1}\bigr)$; then $w_0(\tau_0)=\tfrac2{25}$. If a
covering has $\max_i r_i\ge1-w_0$, then $\sum_ir_i\ge1$, with equality only
for a trivial covering (one $r_i=1$, the others $0$).
\end{lemma}
\begin{proof}
This is the extreme-regime theorem, proved in full in \S[extreme] (the big
plank leaves two corner pieces; the other two planks cover them, reduced to
the two-direction case, Theorem~\ref{thm:twodir}, on each piece, with a
spanning argument when a plank meets both). Reproduced there for general
cyclic $\tau$.
\end{proof}

Write a covering as a point $(l_1,h_1,l_2,h_2,l_3,h_3)\in[0,1]^6=:\mathcal C$.
A covering is \emph{balanced} if $r_i<1-w_0$ for all $i$ and \emph{proper} if
$0<r_i$ for all $i$. Bisection produces axis-aligned boxes $B\subseteq\mathcal
C$; for a box $B$ write $l_i^-\le l_i^+$, $h_i^-\le h_i^+$ for the coordinate
ranges. Define the four \emph{pruning predicates}:
\begin{itemize}
\item[$\mathrm{P1}(B)$:] $\sum_i\max(0,h_i^--l_i^+)>1$.
\item[$\mathrm{P2}(B)$:] $h_i^--l_i^+\ge1-w_0$ for some $i$.
\item[$\mathrm{P4}(B)$:] $l_i^->h_i^+$ for some $i$.
\item[$\mathrm{P3}(B)$:] the \emph{enlarged configuration}
$I_i^+=[l_i^-,h_i^+]$ does not cover $\Delta$.
\end{itemize}

\begin{lemma}[soundness of the prunes]\label{lem:prunes}
If any of $\mathrm{P1}$--$\mathrm{P4}$ holds for a box $B$, then no proper
covering in $B$ has $\sum_ir_i\le1$; more precisely:
\begin{itemize}
\item[\rm(P1)] every configuration in $B$ has $\sum_ir_i\ge\sum_i\max(0,
h_i^--l_i^+)>1$.
\item[\rm(P2)] every configuration in $B$ has $r_i\ge h_i^--l_i^+\ge1-w_0$,
so lies in the extreme regime and satisfies $\sum r\ge1$ by
Lemma~\ref{lem:extreme}.
\item[\rm(P4)] every configuration in $B$ has $l_i>h_i$, i.e.\ plank $i$ is
empty; the covering uses at most the two non-parallel directions $j,k$, so
$\sum r\ge r_j+r_k\ge1$ by Theorem~\ref{thm:twodir}.
\item[\rm(P3)] for every configuration in $B$, $[l_i,h_i]\subseteq I_i^+$, so
$\bigcup_iP_i\subseteq\bigcup_i\{u_i\in I_i^+\}$; if the latter misses a point
of $\Delta$ so does the former --- the configuration is not a covering.
\end{itemize}
\end{lemma}
\begin{proof}
(P1) $r_i=h_i-l_i\ge h_i^--l_i^+$ and $r_i\ge0$. (P2), (P4) as stated,
using $r_i\ge h_i^--l_i^+$ and $l_i\ge l_i^-$, $h_i\le h_i^+$ respectively.
(P3) monotonicity of the covered set in the plank intervals.
\end{proof}

\begin{lemma}[complete exact non-covering test]\label{lem:noncover}
A finite union $\bigcup_i\{u_i\in I_i^+\}$ fails to cover $\Delta$ if and
only if either (a) some edge of $\Delta$ is not covered by the traces of the
$I_i^+$ (a one-dimensional interval computation), or (b) some sign cell
$\{u_i<l_i^+\ \mathrm{or}\ u_i>h_i^+\}_i$ has positive area in $\Delta$.
Both tests are exact and rational.
\end{lemma}
\begin{proof}
An uncovered point $x$ satisfies $u_i(x)\notin I_i^+$ for all $i$, i.e.\ lies
in one sign cell (a relatively open convex polygon in $\Delta$). If $x$ is
interior to $\Delta$ its cell has nonempty interior, hence positive area:
case (b). If $x\in\partial\Delta$ it lies on an edge whose traces miss it:
case (a). Conversely each condition exhibits an uncovered point (a
positive-area cell contains interior points; an edge gap is uncovered).
Areas and interval endpoints are rational functions of $\tau,l^+,h^+$.
\end{proof}

\begin{lemma}[reduction to an empty tree]\label{lem:reduction}
Let $T(\tau_0)$ be the binary tree grown from the root $\mathcal C$ by:
prune a box if any of $\mathrm{P1}$--$\mathrm{P4}$ holds; otherwise bisect its
widest coordinate at the midpoint. Then $T(\tau_0)$ is finite, and: there is a
proper balanced covering of $\Delta$ with $\sum r\le1$ if and only if
$T(\tau_0)$ has a leaf on which no pruning predicate holds.
\end{lemma}
\begin{proof}
Finiteness: a box all of whose coordinates have width $<\eta$ for small enough
$\eta$ satisfies one of the predicates (a proper balanced covering, if it
existed, would have a neighbourhood of coverings by continuity, but every such
box either has $\min\sum r>1$, or its enlarged configuration still fails to
cover --- the quantitative statement is Lemma~\ref{lem:termarg}); so bisection
terminates. Equivalence: by Lemma~\ref{lem:prunes} a pruned box contains no
proper covering with $\sum r\le1$; conversely an unpruned leaf is a box on
which none of the predicates certifies absence, and the search continues, so
leaves of $T$ are exactly the pruned boxes. A proper balanced covering with
$\sum r\le1$ lies in some leaf, which is then unpruned.
\end{proof}

\begin{lemma}[termination as proof]\label{lem:termarg}
If $T(\tau_0)$ has no unpruned leaf, then no proper balanced covering of
$\Delta$ has $\sum r\le1$; hence every proper balanced covering has
$\sum r>1$.
\end{lemma}
\begin{proof}
Immediate from Lemma~\ref{lem:reduction}: an empty tree (all leaves pruned)
leaves no room for such a covering. Strictness: a proper balanced covering
lies in some leaf, pruned by P1, P2, or P3; it is not P4 (proper) and, being a
covering, not P3 (Lemma~\ref{lem:prunes}(P3)); so P1, whose \emph{closed} box
has $\min\sum r>1$, giving $\sum r>1$.
\end{proof}

\begin{theorem}[$C_\Delta(\tau_0)=1$; partially computer-assisted]
\label{thm:sandwich-bang}
For $\tau_0=(\tfrac{13}{25},\tfrac12,\tfrac12)$, every covering of $\Delta$ by
three planks, one parallel to each direction, satisfies $\sum_i\rw\ge1$, with
equality if and only if the covering is trivial (one plank is all of $\Delta$
and the other two are empty). In particular $C_\Delta(\tau_0)=1$: this is the
first non-concurrent triple on the triangle for which the three-direction case
of Conjecture~\ref{conj:bang} is established.
\end{theorem}
\begin{proof}
Partition coverings into: (i) extreme ($\max r_i\ge1-w_0$) ---
Lemma~\ref{lem:extreme}, $\sum r\ge1$, equality trivial; (ii) some plank empty
--- Theorem~\ref{thm:twodir}, $\sum r\ge1$; (iii) proper balanced --- by
Lemmas~\ref{lem:reduction}--\ref{lem:termarg}, $\sum r>1$, provided
$T(\tau_0)$ has no unpruned leaf. \emph{The last is a finite verification in
exact rational arithmetic: the bisection tree $T(\tau_0)$ is empty.} It was
carried out (the searcher of Appendix~\ref{app:cert}) and independently
re-verified from a certificate by a checker that shares no code with the
searcher and re-derives every box and every prune (Appendix~\ref{app:cert});
the emptiness was further corroborated by a second search over a relabelled
triple with a different split ratio, producing a different tree with the same
verdict.

Equality: (i) and (ii) give equality only trivially; among proper two-plank
coverings ($r_k=0$) the minimum width is $\tfrac{31}{25}$ for the pairs
containing $u_1$ and $\tfrac54$ for $(u_2,u_3)$ (both $>1$; exact, ancillary),
so no proper tight two-plank covering exists; and (iii) gives strict
inequality. Hence equality forces a trivial covering.
\end{proof}

\begin{remark}[honest human/machine boundary]\label{rem:boundary}
Every step above is proved by hand except the single assertion ``$T(\tau_0)$
has no unpruned leaf'', which is a finite emptiness check of an explicitly
defined tree. That check is \emph{exact} (rational arithmetic throughout,
independent of hardware and of any floating-point rounding) but not
human-scale ($\approx1.6\times10^6$ tree nodes); it is what makes
Theorem~\ref{thm:sandwich-bang} \emph{partially computer-assisted}. Its
validity rests on the independent checker of Appendix~\ref{app:cert}, not on
trusting the search. This is a single triple, not the family: the general
cyclic case remains open (Problem~\ref{prob:covconst}).
\end{remark}
